\section{Models, certificates, and elementary structure}\label{s1:foundations}

The requirement that all streams leaving a pool have the same
composition is local, but it couples decisions at different products.
Which complexity results apply depends on the precise conventions for
flow contracts, quality specifications, and direct arcs, so we fix them
first. The basic formulations and the rank-one interpretation are
established background. The destination decomposition shows which parts
of profitable-flow detection follow from existing single-product and
multi-commodity arguments; the conic-hull, cyclic, and
facial-specification results give the further structural consequences
used later. None of the sign tests solves the general
objective-threshold problem.

\subsection{Physical models and decision problems}\label{s1:model}

Let $I,L,J$ be disjoint finite sets of inputs, pools, and products.
In the \emph{standard pooling model}, the arc set satisfies
\[
 A\subseteq(I\times L)\cup(L\times J)\cup(I\times J).
\]
An arc in $I\times J$ is a \emph{bypass}. The bypass graph is the
undirected bipartite graph on $I\cup J$ with these arcs as edges.
Its degree, the \emph{bypass degree}, differs from the total degree of
an external node, which also counts incident pool arcs. Write $y_{i\ell},v_{\ell j},z_{ij}$ for
input--pool, pool--product, and bypass flows, respectively. A missing arc
has flow zero. Input $i$ has a rational quality vector
$\lambda_i\in\Q^K$. The quality of pool $\ell$ is $q_\ell\in\R^K$.
Define input withdrawals, pool throughputs, and product deliveries by
\begin{align}
 s_i&=\sum_\ell y_{i\ell}+\sum_jz_{ij},\qquad
 T_\ell=\sum_i y_{i\ell}=\sum_jv_{\ell j},\label{s1:mass}\\
 t_j&=\sum_\ell v_{\ell j}+\sum_i z_{ij}.
\end{align}
The perfect-mixing equations and product quality masses are
\begin{align}
 \sum_i\lambda_{ik}y_{i\ell}&=q_{\ell k}T_\ell
       &&(\ell\in L, k=1,\ldots,K),\label{s1:mixing}\\
 m_{jk}&=\sum_\ell q_{\ell k}v_{\ell j}
                +\sum_i\lambda_{ik}z_{ij}.
\end{align}
Ordinary lower and upper specifications mean
\begin{equation}
 \underline\mu_{jk}t_j\le m_{jk}\le\overline\mu_{jk}t_j.
 \label{s1:quality}
\end{equation}
For the geometric results below we also allow a rational polyhedral
product region $R_j=\{u:A_ju\le b_j\}$, represented in flow variables by
$A_jm_j\le b_jt_j$. This broader specification format is used only when
explicitly stated. A scalar quality with both lower and upper bounds
counts as one attribute. Replacing a lower bound by an upper bound on the
negative of that attribute creates another coordinate in an
upper-bound-only formulation; these are different restrictions.

All flows are nonnegative. Individual arcs and the totals $s_i,T_\ell,t_j$
may have finite rational lower and upper bounds. Initially intersect
each such flow-bound interval with $[0,\infty)$ and reject the instance
if any intersection is empty. In every capacitated
model, these bounds give a finite rational upper bound on every arc flow,
either directly or through finite node capacities. The explicitly
uncapacitated sets introduced below are exempt from this convention.
A bound with equal
endpoints is an \emph{exact flow contract}. An \emph{exact product
contract} fixes both $t_j=d_j$ and $m_j=d_jB_j$ for given demand $d_j$
and quality vector $B_j$. A \emph{common pool bound} constrains $T_\ell$;
it is additional to individual feed and outlet bounds. A theorem that
assumes this bound redundant does not allow an arbitrary restrictive
capacity in its place.

Product quality constraints use their homogeneous mass form even at
zero flow. Since incoming flows are nonnegative, $t_j=0$ forces every
incoming flow to zero, so $m_j=0$ and the constraints impose no composition
requirement. In particular, an exact quality assigned to a zero-demand
product is immaterial. Likewise, $T_\ell=0$ makes every incident pool
flow zero. Such a pool may receive any finite quality without changing
feasibility. Removing an unusable pool requires checking that zero belongs
to its throughput interval and every incident-arc interval; external-node
contracts must remain, since bypasses or other pools can satisfy them.

Costs are rational and linear in actual arc flows. We minimize $c^Tf$,
where $f$ collects those flows, and use profit $-c^Tf$ when convenient.
\emph{Standard economic costs} have the more restrictive form
$\sum_i a_i s_i-\sum_j r_jt_j$. Arbitrary arc costs need not have this
form. Under exact source and product flow contracts, standard economic
cost is constant; arbitrary routing costs need not be constant.

The \emph{threshold problem} asks whether a feasible flow satisfies
$c^Tf\le\tau$ for a rational threshold $\tau$. \emph{Feasibility} omits
this last inequality. With zero lower flow bounds and no positive flow
contracts, the zero flow is feasible. Hardness of a nonzero profit
threshold therefore does not imply hardness of feasibility for that
same class. The \emph{profitable-flow problem} asks whether some feasible
flow has $c^Tf<0$; equivalently, when zero is feasible, whether the optimum
is strictly negative.

We also consider a \emph{generalized pooling network}, with
$A\subseteq(I\cup L)\times(L\cup J)$, which permits pool--pool arcs.
For a flow $f_{uv}$, put $\xi_u=\lambda_u$ at an input and $\xi_u=q_u$
at a pool. The corresponding equations are
\begin{equation}
 F_v:=\sum_u f_{uv}=\sum_w f_{vw},\qquad
 F_vq_v=\sum_u f_{uv}\xi_u \quad(v\in L),\qquad
 m_j=\sum_u f_{uj}\xi_u.\label{s1:generalized}
\end{equation}
An acyclic generalized network has no directed cycles. The cyclic
extension in Section~\ref{s1:cycles} uses the algebraic steady-state
interpretation of \eqref{s1:generalized}, which permits an isolated
positive circulation without source withdrawal. No theorem for that
interpretation assumes an additional source-tracing or initialization
condition. None of the models here includes losses, fixed charges,
nonlinear blending laws, or intermediate quality restrictions unless
explicitly added.

\subsection{Encoding and exact certificates}\label{s1:encoding}

Input size is the binary encoding length of the graph and all rational
data. Polynomial time means bit complexity in this size. When dimensions
are declared fixed, constants and polynomial exponents may depend on
them. Such statements do not imply fixed-parameter tractability with a
uniform exponent, or a practical running-time bound.

An exact real-algebraic answer specifies real roots of rational
polynomials with isolating data and exact expressions for the coordinates.
Whenever we claim a polynomial-time exact algorithm, its output
representation, including degrees, coefficient lengths, and isolating
data, must also have polynomial size. Rational data do not by themselves
imply a rational optimal flow. Conversely, membership in $\NP$ does not
require the certificate to be the flow: a short combinatorial certificate
may be checked by a fixed-dimensional algebraic computation.

For standard and acyclic generalized pooling, active pool qualities are
convex combinations of input qualities. This follows directly by
weighted averaging in standard pooling and by induction in a topological
order in an acyclic network. Inactive pool qualities can therefore be
chosen in the coordinatewise input-quality box. If $I$ is empty, acyclicity
and conservation force all flows to zero, which is checked directly.
When $I$ is nonempty, the finite flow bounds together with this box make
the lifted feasible set closed and bounded. Thus a feasible instance
has an attained linear optimum, and its arc-flow projection is compact.
This argument does not assign source-derived bounds to isolated cyclic
circulations.

The polynomial formulation \eqref{s1:mass}--\eqref{s1:quality} immediately
places the threshold problem in $\ETR$, the class of decision problems
polynomially reducible to existence of a real solution of rational
polynomial equalities and inequalities. This observation is an upper
bound; it supplies no general $\NP$ certificate bound.

A useful exact reduction replaces the attribute count by affine rank.
Suppose $I\ne\varnothing$, and compute rational vectors $\lambda_0$ and
independent columns of $B\in\Q^{K\times r}$ such that
$\lambda_i=\lambda_0+Ba_i$, $a_i\in\Q^r$. Rational Gaussian elimination
computes these data with polynomial encoding length. For standard
pooling, write $q_\ell=\lambda_0+B\eta_\ell$. Equations
\eqref{s1:mass} and \eqref{s1:mixing} are then equivalent to
\begin{equation}
 \sum_i a_i y_{i\ell}=\eta_\ell T_\ell.\label{s1:rank-coordinate}
\end{equation}
Indeed, substitution cancels $\lambda_0T_\ell$, and independence of the
columns of $B$ proves the converse. Every active pool is represented;
an inactive pool can be assigned arbitrary coordinates in the box spanned
coordinatewise by the $a_i$. Substitution into every original product
inequality preserves all attributes, including when their number grows.
If $r=0$, all input qualities coincide and quality feasibility reduces
to deletion of incompatible positive product flows and linear constraints.
Affine quality rank and the number of distinct quality vectors are
separate parameters.

Readers interested mainly in algebraic completeness can continue with
Section~\ref{s2:algebraic}; the remaining subsections of this section
are independent of those constructions.

\subsection{The rank-one flow block}\label{s1:rank-one}

For one standard pool, let $y_i$ and $v_j$ be its incoming and outgoing
flows, with common total $T$. When $T>0$, the amount of input $i$
allocated to product $j$ through that pool is
\begin{equation}
 W_{ij}=\frac{y_iv_j}{T}.\label{s1:rank-one-form}
\end{equation}
Thus $W$ is nonnegative, has rank one, and has row and column sums $y,v$.
When $T=0$, nonnegativity gives $y=v=0$, and we set $W=0$.
Conversely, every nonzero nonnegative rank-one matrix has
\eqref{s1:rank-one-form} with its own margins. To see this, write
$W=ab^T$ with $a,b\ge0$; such factors follow by taking a positive column
and the proportional nonnegative column multipliers. Its margins are
$y=a\sum_jb_j$, $v=b\sum_i a_i$, and
$T=(\sum_i a_i)(\sum_jb_j)$, proving the identity.
The mass of attribute $k$ delivered along the pool's outlet to $j$ is
$\sum_i\lambda_{ik}W_{ij}$. This explains the role of rank-one matrix
sets with row and column bounds in pooling formulations
\citep{dey2020-convexifications-of-rank-one-based}.

This correspondence requires care with objectives. Physical costs on
feed and outlet arcs induce the matrix coefficients $C_{ij}=a_i+b_j$.
An arbitrary cost matrix on $W$ is a broader optimization model.
In particular, hardness of arbitrary linear optimization over a
rank-one block alone does not prove hardness of a single physical pool
with only feed and outlet costs.

\subsection{Single-product decomposition and profitable flow}\label{s1:triviality}

Throughout this subsection the generalized network is acyclic, flow
lower bounds are zero, and all other flow restrictions are nonnegative
upper capacities. Only final products have quality restrictions. Let
$P$ be the set of feasible arc flows and $z^*=\min_{f\in P}c^Tf\le0$.
The destination-disaggregation construction is established pooling
machinery \citep{boland2016-new-multi-commodity-flow-formulations}.
The standard-pooling one-product
approximation principle is due to
\citet{dey2015-analysis-of-milp-techniques-for}; the sign question appears
explicitly in \citet[Remark~2.2]{gupte2017-relaxations-and-discretizations-for-the}.
\citet[Proposition~4 and Theorem~2]{xiang2026-dive-and-fix} resolve that
question by single-terminal LPs in a generalized terminal-proportion
formulation, including pool--pool arcs, so polynomial zero-optimum
detection is established. We give the full argument because it works
directly with physical flows and supplies the decomposition used for
the shortest-path and conic-hull statements.

\begin{lemma}[Destination decomposition]\label{s1:destination}
Every feasible physical flow has a decomposition
$f=\sum_{j\in J}f^j$ into feasible flows, each serving only product $j$,
with $0\le f^j\le f$ coordinatewise.
\end{lemma}
\begin{proof}
Fix feasible pool qualities. At products set
$h_j(j')=1$ if $j=j'$ and zero otherwise. In reverse topological order,
set, for each positive-throughput pool $v$,
\[
 h_j(v)=F_v^{-1}\sum_wf_{vw}h_j(w).
\]
Set all $h_j(v)$ to zero at zero-throughput pools. Conservation and
acyclicity imply that every positive-flow path can be continued to a
product; hence reverse induction gives $0\le h_j(v)\le1$ and
$\sum_jh_j(v)=1$ at each active pool. Define
$f^j_{uv}=f_{uv}h_j(v)$, using the head of the arc.
All incoming flows of a pool $v$ receive the same multiplier $h_j(v)$.
Their sum is $F_vh_j(v)$; the sum of outgoing component flows has this
same value by the recursion. Their incoming quality mass is the original
quality mass multiplied by $h_j(v)$, so the original $q_v$ still satisfies
\eqref{s1:generalized}. At product $j$ the original inflow is unchanged;
other products receive zero. Coordinatewise domination preserves all
upper capacities, including node throughputs. Finally, the head of any
positive-flow arc is either an active pool or a product, where the
fractions sum to one. Summing the components therefore gives $f$.
\end{proof}

For each $j$, form a rational LP by retaining all ordinary flow balances
and upper capacities, forcing delivery to other products to zero, and
replacing all mixing equations by
\begin{equation}
 \underline\mu_{jk}t_j\le\sum_i\lambda_{ik}s_i
                   \le\overline\mu_{jk}t_j\quad(k=1,\ldots,K).
 \label{s1:single-lp}
\end{equation}
For polyhedral product regions the analogous row is
$A_j\sum_i\lambda_is_i\le b_jt_j$.

\begin{lemma}[Single-product projection]\label{s1:single}
This LP describes exactly the arc flows of the physical single-product
problem.
\end{lemma}
\begin{proof}
Summing all physical pool quality balances cancels internal quality
flows, so the quality mass reaching the only served product is
$\sum_i\lambda_is_i$. This proves necessity. Conversely, take an LP
flow and process pools in topological order. Give every active pool the
weighted average of its incoming qualities and assign arbitrary finite
qualities at inactive pools. All pool equations hold. Summing them again
identifies product quality mass with the expression in
\eqref{s1:single-lp}, proving its specification. Zero delivery forces
zero flow throughout an acyclic conserved network, so that case also
satisfies all constraints.
\end{proof}

\begin{theorem}[Polynomial sign test; physical-flow version]\label{s1:sign}
Let $z_j$ be the minimum cost in the single-product LP. Then
$z^*=0$ if and only if every $z_j=0$. These LPs also construct a
profitable physical flow whenever one exists. If $m=|J|\ge1$ and
$z_{\rm best}=\min_jz_j$, then
\begin{equation}
 z^*\le z_{\rm best}\le z^*/m\le0.\label{s1:approx}
\end{equation}
\end{theorem}
\begin{proof}
Every single-product LP flow is physical by Lemma~\ref{s1:single}.
Conversely, Lemma~\ref{s1:destination} decomposes an optimal physical
flow into $m$ feasible components whose costs sum to $z^*$.
At least one component has cost at most $z^*/m$, and its LP optimum is
no larger. This proves \eqref{s1:approx} and the sign equivalence.
Rational LP optimization and exact comparison with zero take polynomial
bit time. Reconstructing qualities as in Lemma~\ref{s1:single} is exact.
For rational flow data this reconstruction solves a rational triangular
linear system on active pools, so determinant bounds give polynomial
encoding length. If $J$ is empty, acyclicity forces zero flow directly.
\end{proof}

For $m=|J|\ge1$, there is also a single destination-flow relaxation. Give every product
its own nonnegative conserved single-product flow $f^j$ satisfying
\eqref{s1:single-lp}, and impose all capacities on $\sum_j f^j$.
Let $z_{\rm rel}$ be its minimum cost. Destination decomposition makes
this a relaxation. Every component of a feasible relaxed solution is
physical and respects the capacities by nonnegativity. Selecting the
cheapest component of an optimal relaxed solution proves
\[
 z_{\rm rel}\le z^*\le z_{\rm best}\le z_{\rm rel}/m\le0.
\]
In particular, this relaxation has the same zero-optimality status as
pooling. These inequalities concern cost signs and product-count profit
approximation; they are not exact optimization formulas for pooling.

\begin{theorem}[Shortest-path sign criterion]\label{s1:paths}
Delete zero-capacity arcs and nodes with their incident arcs. Let
$\delta_{ij}$ be the minimum path cost from input $i$ to product $j$
on the remaining graph, and let $I_j$ contain the inputs that reach $j$.
Define $\beta_j$ as the optimum of
\begin{align*}
 \min\ &\sum_{i\in I_j}\delta_{ij}\gamma_i\\
 \text{subject to }&\gamma\ge0,\quad\sum_{i\in I_j}\gamma_i=1,\quad
 \underline\mu_{jk}\le\sum_{i\in I_j}\lambda_{ik}\gamma_i
                                  \le\overline\mu_{jk}\quad(k=1,\ldots,K),
\end{align*}
with value $+\infty$ for an infeasible LP. Then $z^*<0$ if and only if
some $\beta_j<0$. A profitable instance has a rational polynomial-size
witness obtained by routing at most $K+1$ shortest input--product paths
and scaling all flows by one positive factor.
\end{theorem}
\begin{proof}
A negative-cost physical flow has a negative-cost single-product
component. Decompose its ordinary conserved flow into input--product
paths and divide source withdrawals by its positive total delivery $t$.
The resulting weights $\gamma_i=s_i/t$ meet the quality bounds by
aggregate conservation. Actual cost per unit is at least
$\sum_i\delta_{ij}\gamma_i$, so $\beta_j<0$.
Conversely, route a negative-cost feasible mixture on one shortest path
per supported input. Every retained capacity is strictly positive, so
a sufficiently small positive common scale satisfies all capacities
and preserves cost sign and quality ratios. The scaled conserved
single-product flow is physical by Lemma~\ref{s1:single}.

The mixture LP is compact. Choose an optimal vertex and let $s$ be its
support size. Its active quality rows have rank at most $K$, and
normalization adds one row. If $s>K+1$, a nonzero direction on its
support preserves every active row. Small perturbations in both signs
preserve inactive inequalities and nonnegativity, contradicting
extremality. Thus $s\le K+1$. Shortest paths in an acyclic graph have
polynomial length and rational polynomial-size costs. The mixture
vertex has polynomial rational encoding. For each capacity row $r$, let
$U_r$ be its upper bound and $a_r$ its load in the unscaled flow. Choose
the scale as $\min(\{1\}\cup\{U_r/a_r:a_r>0\})$; zero-load rows impose
no restriction. This positive rational scale has polynomial encoding.
Quality reconstruction has
the bit bound established in Theorem~\ref{s1:sign}.
\end{proof}

The support bound uses $K$ coordinate quality bounds. For general
polyhedral specifications it becomes one plus the rank of the active
quality rows restricted to the input simplex. Positive capacity
magnitudes affect attainable profit but affect its sign only through
the support remaining after zero-capacity deletion.

\subsection{An exact uncapacitated conic hull}\label{s1:conic}

Continue on the positive-capacity support. Let $S$ be the physical flow
set after removing all upper capacities, and let $C_j$ be its
single-product LP cone. We use $\cone(X)$ for finite nonnegative linear
combinations of points of $X$; it is not assumed closed by definition.

\begin{theorem}\label{s1:cone-theorem}
For the acyclic model of Section~\ref{s1:triviality},
\begin{equation}
 \conv(S)=\cone(S)=\sum_{j\in J}C_j=\cone(P).\label{s1:cone-equality}
\end{equation}
The sum is a Minkowski sum and has a polynomial-size linear extended
formulation. In particular, this cone is polyhedral and closed.
\end{theorem}
\begin{proof}
Destination decomposition proves $S\subseteq\sum_jC_j$.
Every $C_j\subseteq S$ by Lemma~\ref{s1:single}, so
$\sum_jC_j\subseteq\cone(S)$. Because $S$ contains zero and is closed
under nonnegative scaling, any finite conic combination can be written
as a convex combination of scaled elements of $S$; hence
$\cone(S)=\conv(S)$. The Minkowski sum is a convex cone, proving the
first three equalities. Since $P\subseteq S$, one inclusion for the
last equality is immediate. Conversely every finite $f\in S$ can be
scaled by a positive scalar into $P$, since all retained capacities
are positive. Thus $S\subseteq\cone(P)$. The formulation keeps a copy
of the $O(|A|)$ flow variables for each product and the equation
$f=\sum_jf^j$. It is a projection of a polyhedral cone and hence is
polyhedral. With no products every cone in question is $\{0\}$.
\end{proof}

Consequently all homogeneous linear inequalities valid for $P$ are
captured by \eqref{s1:cone-equality}. The profitable-flow test asks whether
$c$ lies outside its dual cone
$\{c:c^Tf\ge0\text{ for every }f\in S\}$; membership is equivalent
to the absence of profitable flow.

Capacities cannot generally be reimposed after convexification to obtain
$\conv(P)$. For example, take input qualities zero and one, a single
pool of capacity two, and two products requiring qualities exactly zero
and exactly one. Allow only feed and outlet arcs and give each input,
product, and arc capacity one. A physical flow serves at most one
product, so total delivery is at most one throughout $\conv(P)$.
The sum of the two unit pure single-product flows belongs to
$\conv(S)$ and satisfies every stated capacity, but delivers two units.

\subsection{Algebraic recirculation}\label{s1:cycles}

The sign and conic-hull conclusions extend to cyclic generalized
networks under the steady-state semantics of
Section~\ref{s1:model}. All assumptions on zero lower bounds and
absence of intermediate quality restrictions remain in force.
The generalized sign test of \citet{xiang2026-dive-and-fix} already
allows pool--pool arcs. Here we specify how the physical equations
treat a singular case that a terminal-indexed formulation must handle
separately: an isolated positive circulation, including one with no
path to any product. This also identifies its contribution to the
uncapacitated cone.

\begin{lemma}[Cyclic single-product reconstruction]\label{s1:cyclic-single}
Every ordinary nonnegative conserved flow serving at most one product
admits pool qualities satisfying the mixing equations. Its delivered
quality mass is $\sum_i\lambda_is_i$. With rational flow data,
qualities of polynomial rational encoding can be constructed.
\end{lemma}
\begin{proof}
Consider the strongly connected components of the positive pool-flow
support, ignoring inactive pools. Summing mass balances over any
component shows that its total external inflow equals its total external
outflow. A component with no external inflow therefore has no external
outflow and is an isolated circulation. Assign a constant quality to
all its pools. Process the other components in topological order.
For one coordinate, the equations on a component $D$ are
\[
 F_vq_v-\sum_{u\in D}f_{uv}q_u=b_v,\qquad
 b_v=\sum_{u\notin D}f_{uv}\xi_u.
\]
After division by $F_v>0$, the internal matrix $Q$ is nonnegative and
has row sums at most one, with at least one deficient row. Strong
connectivity implies that every state of the corresponding backward
chain has a positive-probability path to a deficient row and then out
of the component. Finiteness gives an integer $h$ and $\varepsilon>0$
such that every state exits within $h$ steps with probability at least
$\varepsilon$. Thus $\|Q^{rh}\|_\infty\le(1-\varepsilon)^r$, proving
that $I-Q$ is invertible with inverse $\sum_{a\ge0}Q^a$.
The solution is a weighted average of the already known boundary
qualities. This constructs all active pool qualities; inactive ones
are arbitrary. Summing the quality balances cancels all internal mass,
including cycles, and proves the aggregate identity.
For rational flows, the equations on all nonisolated active pools form
one nonsingular rational linear system. Cramer's rule gives polynomial
encoding length; isolated components can use quality zero.
\end{proof}

\begin{theorem}[Cyclic sign and hull]\label{s1:cyclic-sign}
Let $C_0$ be the nonnegative circulation cone on the pools and $C_j$
the uncapacitated single-product cone, now allowing cycles. Then
\[
 \conv(S)=\cone(S)=C_0+\sum_{j\in J}C_j=\cone(P).
\]
Profitable flow can be detected in polynomial time. On the
positive-capacity support it exists precisely when there is a
negative-cost directed cycle, or, in the absence of such a cycle,
one of the shortest-path mixture LPs in Theorem~\ref{s1:paths} has
negative optimum.
\end{theorem}
\begin{proof}
For any feasible flow, a terminal strongly connected component of its
positive pool support with no positive arc to a product has no positive
external inflow by conservation. It is therefore an isolated circulation.
Remove all such components and call their flow $f^0$. Every remaining
active pool reaches a product along positive arcs. The forward routing
chain with transition probabilities $f_{vw}/F_v$ is absorbed at products
with probability one: as in Lemma~\ref{s1:cyclic-single}, there is a
uniform positive chance of absorption within a fixed number of steps.
Its absorption probabilities satisfy
\[
 F_vh_j(v)=\sum_wf_{vw}h_j(w),\qquad\sum_jh_j(v)=1.
\]
Set $h_j(v)=0$ on the removed circulation components and inactive pools,
and set $h_j(j')=1$ if $j=j'$ and zero otherwise at products.
Using $f^j_{uv}=f_{uv}h_j(v)$ gives, by the same incoming-mass
calculation as Lemma~\ref{s1:destination}, physical dominated components
and $f=f^0+\sum_jf^j$. This proves one containment for the hull.
Every nonnegative circulation is physical, by choosing a common constant
quality on all its pools, and each $C_j$ is physical by
Lemma~\ref{s1:cyclic-single}. Scaling and convexification then prove all
hull equalities exactly as before.

A negative directed cycle is a feasible profitable circulation after
small positive scaling. If there are no negative cycles, $f^0$ has
nonnegative cost by cycle decomposition. A profitable flow therefore
has a negative single-product component. Its ordinary path-and-cycle
decomposition, followed by removal of nonnegative-cost cycles, yields
a negative path mixture with unchanged source withdrawals and delivery.
The aggregate quality identity then proves a negative shortest-path
mixture value. Conversely, route any such mixture on shortest paths
and scale to capacities. Lemma~\ref{s1:cyclic-single} supplies feasible
qualities even when the union of the paths contains cycles. Negative-cycle
detection and shortest-path algorithms with signed arc costs, followed
by rational LPs, take polynomial time.
\end{proof}

If $J\ne\varnothing$, $C_0$ is contained in every $C_j$, so its extra
summand is redundant. The capacitated single-product LP family also
detects negative circulations, since each LP admits them. If $J$ is
empty, the sign problem is an ordinary circulation problem.

For $m=|J|\ge1$, the approximation and relaxation inequalities of
Section~\ref{s1:triviality} hold in the cyclic model as well. In the
decomposition used above, merge $f^0$ into one designated product
component. Its circulation support is disjoint from the other components'
positive support, so the merged component is physical and is still
coordinatewise dominated by $f$. There are now exactly $m$ feasible
single-product components summing to $f$. Applying the averaging argument
to an optimal physical flow proves \eqref{s1:approx}. The same
decomposition makes the destination-flow LP a relaxation, and each of
its components is physical by Lemma~\ref{s1:cyclic-single}. Selecting
its cheapest component therefore gives
$z_{\rm rel}\le z^*\le z_{\rm best}\le z_{\rm rel}/m\le0$.
An optimum exists here as well: nonisolated pool qualities lie in the
input-quality hull, and isolated circulations can all use quality zero.
These choices put every feasible flow in the projection of one bounded
closed lift, using the box containing zero and all input qualities
(the singleton zero box when there are no inputs).

The distinction between active support and the original graph is
essential. On the arcs $i\to a$, $a\to b$, $b\to a$, $b\to j$, let
$f_{ab}=f_{ba}=1$ and the other flows be zero. Both pool qualities may
equal any common constant, although every vertex lies on an
input--product path in the original graph. The mixing equations have
a singular cycle block. More quantitatively, take
$f_{ia}=\varepsilon$, $f_{ab}=1$, $f_{ba}=1-\varepsilon$,
$f_{bj}=\varepsilon$, with $0<\varepsilon<1$ and input quality $\lambda$.
The pool equations are
\[
 \begin{pmatrix}1&-(1-\varepsilon)\\-1&1\end{pmatrix}
 \begin{pmatrix}q_a\\q_b\end{pmatrix}
 =\begin{pmatrix}\varepsilon\lambda\\0\end{pmatrix}.
\]
Their determinant is $\varepsilon$ and the solution is
$q_a=q_b=\lambda$. Assigning both qualities $\lambda+1$ produces residual
$(\varepsilon,0)^T$ despite unit concentration error. Thus an absolute
mixing residual alone does not control concentration error near a
closed circulation. The treatment of cyclic formulation variables is
part of the existing generalized-pooling literature
\citep[Section~4.3]{boland2016-new-multi-commodity-flow-formulations}; the arguments
here explicitly separate the singular circulation case.

\subsection{Facial product specifications and integral flow}\label{s1:faces}

Return to standard pooling, permitting arbitrary bypasses, rational
polyhedral product regions, and lower as well as upper flow bounds.
Assume a nonempty input set and put
$C=\conv\{\lambda_i:i\in I\}$ and $F_j=C\cap R_j$.
We count $\varnothing$ and $C$ as faces. A face $F$ of $C$ has the
property that a strict convex combination of points of $C$ belongs to
$F$ only if every positively weighted point belongs to $F$.
For example, an upper quality bound zero with nonnegative input
qualities selects a face: a positive-quality feed cannot contribute to
such a product. This exclusion already appears in pooling hardness
reductions, including \citet[Theorem~4]
{haugland2016-the-computational-complexity-of-the}. The result below
identifies exactly when the same reasoning guarantees an integral
optimum for every network and every linear objective on a given set of
input qualities.

\begin{theorem}[Universal flow integrality]\label{s1:facial}
If every $F_j$ is a face of $C$, the projected feasible flow set is a
finite union of network flow polytopes. With integer flow bounds, these
polytopes are integral; consequently every feasible instance has an
integral optimal flow for every rational linear objective.
Conversely, for any nonempty $F=C\cap R$ that is not a face, there is
a one-pool, one-product network on the same input qualities with all
arc and vertex upper bounds one, zero lower bounds, and costs in
$\{0,-1\}$ for which no integral flow is optimal.
\end{theorem}
\begin{proof}
If product $j$ receives positive flow, its blend lies in $F_j$.
The face property forces every contributing pool quality and every
contributing direct input quality into $F_j$. Applying it again to
an active contributing pool forces $\lambda_i\in F_j$ whenever both
$y_{i\ell}>0$ and $v_{\ell j}>0$. Conversely, if this compatibility
holds for every such pair and every positive direct arc, weighted
averaging puts every delivered blend in $F_j$. If $F_j$ is empty,
these rules permit no positive delivery to $j$.

Fix an allowed arc set containing no incompatible direct arc or feed--outlet
pair. All ordinary network flows restricted to that set are quality
feasible by the preceding argument. Taking the finite union over
compatible allowed sets gives exactly the pooling projection. Each
branch has only conservation, arc bounds, and node-throughput bounds.
Split nodes to represent the latter by bounded arcs and add a return
arc from a common terminal to a common source to express the whole
problem as a bounded circulation. A finite upper bound for the return
arc is the sum of source upper bounds when these are supplied; otherwise
the sum of the effective arc upper bounds suffices.

For completeness, a directed incidence matrix is totally unimodular.
In any square submatrix, a column with at most one nonzero entry permits
induction by determinant expansion; if no such column exists, every
column has one $+1$ and one $-1$, so the rows are dependent. Thus every
square subdeterminant is $0$ or $\pm1$. Adding unit rows for bounds
preserves this property by the same expansion. A vertex basis with
integer right-hand side therefore has integer solution by Cramer's rule.
This proves branch integrality, also after projecting auxiliary arcs.
Boundedness and finiteness of the union give an attained integral optimum
whenever at least one branch is nonempty.

For necessity, nonfaciality gives $x,y\in C$, $0<\theta<1$, with
$q=\theta x+(1-\theta)y\in F$ and, say, $x\notin F$. Expand $x$ and
$y$ into input mixtures. Some input outside $F$ has positive weight
in the representation of $x$, since otherwise convexity would put $x$
in $F$. Hence the representation of $q$ uses positive total mass from
inputs outside $F$. Connect all inputs to a single unit-capacity pool
and that pool to a unit-capacity product with region $R$. Charge $-1$
on feeds from inputs outside $F$ and zero on other arcs. The displayed
unit mixture is feasible with negative cost. An integral feasible flow
has throughput zero or one; in the latter case it comes entirely from
one input whose quality must lie in $F$. Every integral feasible flow
therefore has cost zero. The underlying single-product flow constraints
are rational linear constraints, so a rational negative-cost mixture
also exists, for example at an optimal LP vertex.
\end{proof}

The converse is a necessity statement for this universal guarantee,
quantified over networks and objectives, not a hardness theorem. It does
not imply that every given network with a nonfacial product has a
fractional optimum, or that a nonfacial class is computationally hard. The facial
class itself will contain hard support-selection problems. Integer flow
also does not imply pure input selection when a pool can process several
integer units.

\begin{proposition}[Recognition and bounded-pool algorithm]\label{s1:face-alg}
Faciality of all product specifications is recognizable in polynomial
bit time. For fixed pool count $p$ and fixed affine input-quality
dimension $d$, standard pooling with facial specifications admits exact
polynomial-time optimization, including arbitrary bypasses and rational
flow bounds.
\end{proposition}
\begin{proof}
For product $j$, call input $i$ forbidden if $\lambda_i\notin R_j$.
Maximize the total weight of forbidden inputs over
\[
 w\ge0,\quad\sum_iw_i=1,\quad\sum_i\lambda_iw_i\in R_j.
\]
Infeasibility means $F_j=\varnothing$. If feasible, the maximum is zero
for a face by its defining property. Conversely, a zero maximum implies
that every feasible mixture uses only allowed inputs. If an interior
point of a segment in $C$ belongs to $F_j$, expand both endpoints as
input mixtures. A forbidden input in either expansion would yield
positive objective in the LP. Both endpoints therefore are mixtures of
allowed inputs and belong to $F_j$, proving faciality. All these tests
are rational LPs of polynomial size.

Enumerate the nonempty faces $H$ of $C$. For each tuple $(H_\ell)_{\ell\in L}$,
retain feed $(i,\ell)$ only if $\lambda_i\in H_\ell$, outlet $(\ell,j)$
only if $H_\ell\subseteq F_j$, and bypass $(i,j)$ only if
$\lambda_i\in F_j$. Solve the resulting bounded min-cost network
problem. Every resulting flow is physical by convexity. Conversely,
given a feasible physical flow, choose for each active pool the smallest
face of $C$ containing its quality. The face property retains all its
positive feeds; every receiving product face contains the chosen face,
so all its positive outlets are retained. An inactive pool can choose
any nonempty face. If an omitted arc has a positive lower bound, that
branch is infeasible; it must not be removed without this check.

To bound enumeration, work in the rational affine hull of $C$. If
$d>0$, every facet contains $d$ affinely independent input points.
Enumerate all such subsets and test their supporting hyperplanes, giving
at most $|I|^d$ facet candidates. A proper nonempty face is obtained by
intersecting at most $d$ facet hyperplanes with $C$: take a basis of
its active facet normals. The other active hyperplanes follow because
their normals are linear combinations and all equalities agree at a
point of the face. Enumerating at most $d$ facets gives at most
$1+|I|^{d^2}$ face descriptions, including duplicates and $C$ itself.
For $d=0$, $C$ is the only nonempty face. Membership and containment can
be tested on the input points generating each face. Thus a polynomial
number of tuples is inspected for fixed $p,d$, with polynomial-size
rational data and polynomial-time network solves.
\end{proof}

A particularly compact special case has all input attributes in $[0,1]$
and product quality restrictions consisting only of coordinate upper
bounds in $\{0,1\}$. No additional nonredundant quality restrictions
are imposed; lower flow bounds remain permitted. Put
\[
 D_{\ell k}=\sum_{i:\lambda_{ik}>0}y_{i\ell},\qquad
 S_{\ell k}=\sum_{j:\overline\mu_{jk}=0}v_{\ell j}.
\]
Quality feasibility is equivalent to deletion of bypasses with
$\lambda_{ik}>0$ and $\overline\mu_{jk}=0$ for some $k$, together with
\begin{equation}
 D_{\ell k}=0\quad\text{or}\quad S_{\ell k}=0
             \qquad(\ell\in L, k=1,\ldots,K).\label{s1:endpoint-or}
\end{equation}
Indeed, an upper bound zero is a sum of nonnegative quality masses
bounded by zero. Every positive contributing pool must therefore have
quality zero in that coordinate, which excludes every positive-quality
feed. Conversely, the disjunction ensures this exclusion for every
pool serving a strict product; upper bounds one accept every mixture.
This proves both directions, including zero flow.

Choosing one side of each disjunction gives at most $2^{pK}$ ordinary
network problems. An exact MILP uses one binary $b_{\ell k}$ per pair
and, for a finite throughput upper bound $U_\ell$, the rows
\[
 D_{\ell k}\le U_\ell(1-b_{\ell k}),\qquad
 S_{\ell k}\le U_\ell b_{\ell k}.
\]
The formulation is exact, with no quality discretization. With integer
flow bounds, an integral flow is a polynomial-size threshold certificate;
\eqref{s1:endpoint-or} checks its quality feasibility directly. With
rational bounds, a branch LP vertex provides a rational certificate
instead. Thus this endpoint-specification threshold class belongs to
$\NP$, while the number of support branches can still grow exponentially.
